\documentclass[11pt]{article}
\usepackage[T1]{fontenc}
\usepackage{amsmath,amssymb,amsthm}
\usepackage{hyperref}
\usepackage{enumitem}

\newtheorem{theorem}{Theorem}
\newtheorem{lemma}{Lemma}
\newtheorem{proposition}{Proposition}
\newtheorem{corollary}{Corollary}
\theoremstyle{definition}
\newtheorem{definition}{Definition}
\newtheorem{assumption}{Assumption}
\newtheorem{remark}{Remark}

\title{Proof of Partial Correctness and Safety Properties of \texttt{soc-agent}}
\author{Codex}
\date{March 27, 2026}

\begin{document}
\maketitle

\section{Statement of Scope}
This document proves the \emph{partial correctness}, \emph{logical consistency}, and
\emph{fail-closed safety properties} of the implemented \texttt{soc-agent} control plane.
It does \emph{not} prove universal empirical detection accuracy against all real-world
security incidents, because such a claim would require omniscient ground truth over an
open world with external APIs, incomplete evidence, and non-deterministic model behavior.

\begin{definition}[Control-plane correctness]
For this document, \texttt{soc-agent} is \emph{control-plane correct} if the following hold:
\begin{enumerate}[label=(\arabic*)]
\item the investigation task graph actually executed is a deterministic function of the normalized alert;
\item the scheduler never executes a task before its dependencies are satisfied;
\item retries, timeouts, and early-stop behavior are bounded and fail closed;
\item adapter execution is never attempted on the remediation path unless the explicit policy conditions for execution are satisfied;
\item persisted memory, correlation, and replay operations preserve the stored incident semantics;
\item the stated invariants compose without contradiction in \texttt{run\_investigation}.
\end{enumerate}
\end{definition}

\section{Assumptions}
\begin{assumption}[Runtime semantics]
Python, SQLite/Postgres transactions, JSON parsing, and the filesystem behave according to
their documented semantics.
\end{assumption}

\begin{assumption}[Adapter contract]
External model providers and integration adapters either return values conforming to the
interfaces used by \texttt{soc-agent} or raise exceptions.
\end{assumption}

\begin{assumption}[Deduplicated entities on the application path]
On the main application path, the entity lists written into the memory store are deduplicated.
This assumption is satisfied by the implemented path
\[
\texttt{extract\_entities\_from\_graph} \rightarrow \texttt{MemoryStore.write\_memory},
\]
because the extractor accumulates entities in sets and emits sorted unique lists.
\end{assumption}

\begin{remark}
Assumption 3 is needed only for the exact baseline-count monotonicity claim. The main safety
claims do not depend on it.
\end{remark}

\section{Boundary of What Cannot Be Proved}
\begin{proposition}[Universal detection accuracy is not provable here]
No proof from the present source tree can establish that \texttt{soc-agent} is always
factually accurate about whether an arbitrary real-world alert is malicious, benign, or
complete in evidentiary coverage.
\end{proposition}

\begin{proof}
The system depends on external evidence sources, operator configuration, and language-model
outputs. Ground-truth maliciousness is not derivable from the code alone. Therefore the
strong claim ``the system is universally accurate'' is not a theorem of this implementation.
\end{proof}

\begin{remark}
This limitation does not weaken the formal proof below. It means only that the provable
statement must be about implemented semantics and safety, not omniscient real-world truth.
\end{remark}

\section{Planner and Commander}
\begin{lemma}[Planner totality and determinism]
For any input alert object, \texttt{Planner.build\_plan(alert)} returns a finite plan whose
\texttt{plan\_id}, alert type, objective, task identities, dependencies, retry bounds, and
timeouts are deterministic functions of the alert fields used by the planner.
\end{lemma}

\begin{proof}
The planner normalizes only the alert type and severity, computes a UUIDv5 from a fixed
namespace string and alert metadata, chooses one of finitely many hard-coded task templates,
and materializes tasks from that template. No random source, network input, clock value, or
mutable global state participates. Therefore the output is deterministic. The plan is finite
because every branch returns a list of exactly five task specifications.
\end{proof}

\begin{lemma}[Planner-generated dependency graph is acyclic]
Every plan emitted by \texttt{Planner.build\_plan} is a directed acyclic graph.
\end{lemma}

\begin{proof}
Let the canonical order of task roles be
\[
(\texttt{recon}, \texttt{threat\_intel}, \texttt{forensics}, \texttt{remediation}, \texttt{reporter}).
\]
In every planner branch:
\begin{enumerate}[label=(\alph*)]
\item \texttt{recon} has no dependencies;
\item \texttt{threat\_intel} depends only on \texttt{recon};
\item \texttt{forensics} depends only on \texttt{recon};
\item \texttt{remediation} depends only on \texttt{threat\_intel} and \texttt{forensics};
\item \texttt{reporter} depends only on earlier tasks.
\end{enumerate}
Hence every edge points from an earlier position in the above order to a later one. A directed
cycle would require some edge to point backwards in that order, which never occurs. Therefore
the graph is acyclic.
\end{proof}

\begin{lemma}[Commander does not let the LLM choose the task graph]
The actual executed task graph is independent of the commander's LLM output.
\end{lemma}

\begin{proof}
Inside \texttt{Commander.investigate}, the LLM response is parsed into a local variable
\texttt{plan} used only for logging the objective text. The actual execution path is
\[
\texttt{Commander.\_run\_plan(alert, agents\_map)} \rightarrow
\texttt{self.planner.build\_plan(alert)}.
\]
Therefore the scheduled plan topology comes from the deterministic planner, not from the LLM.
If the LLM call fails or returns invalid JSON, the commander falls back to a fixed local
dictionary for logging and still invokes the deterministic planner. Hence the executed DAG is
LLM-independent.
\end{proof}

\section{Scheduler}
\begin{lemma}[Dependency safety]
The scheduler never starts a task unless all of its dependencies are either completed or skipped.
\end{lemma}

\begin{proof}
At each iteration the ready set is defined exactly by
\[
\texttt{all(dep in completed or dep in skipped for dep in task.dependencies)}.
\]
Only tasks in that ready set are passed to \texttt{asyncio.gather} for execution. Therefore a
task cannot start before every dependency is marked completed or skipped.
\end{proof}

\begin{lemma}[Blocked propagation]
If any dependency of a remaining task has status failed or blocked, then that task is marked
blocked and is not executed.
\end{lemma}

\begin{proof}
At the start of every loop, \texttt{\_block\_unreachable\_tasks} scans remaining tasks and checks
whether any dependency belongs to the failed or blocked sets. Every such task is immediately
assigned status \texttt{blocked}, written into \texttt{task\_results}, added to the blocked set,
and removed from the remaining set. Since it is removed before the ready-set computation, it
cannot be executed afterward.
\end{proof}

\begin{lemma}[Bounded retries and timeout fail-closed]
For each task with retry bound \(r\), the scheduler invokes the runner at most \(r+1\) times.
If the runner exceeds its timeout or raises an exception on all attempts, the task finishes in
status \texttt{failed}.
\end{lemma}

\begin{proof}
The scheduler increments the local counter \texttt{attempts} once per loop of
\texttt{\_run\_task}. After each failed attempt it checks whether \texttt{attempts > max\_retries}.
The first time this inequality holds, it returns a failed \texttt{TaskExecutionResult}. Hence
the number of attempts is at most \(r+1\). Timeout exceptions are caught and converted to a
failure string; arbitrary exceptions are also caught and stored as the last error. Therefore
timeout and repeated exception cases fail closed.
\end{proof}

\begin{lemma}[Early stop is optional-only]
Early stop can skip only optional tasks, and only after every mandatory non-reporter task has
completed with average confidence at least the configured threshold, with no failed or blocked
tasks present.
\end{lemma}

\begin{proof}
The predicate \texttt{\_should\_early\_stop} first rejects early stop if any non-optional task
other than \texttt{reporter} is not completed. It then computes the mandatory confidence and
rejects early stop when the confidence is undefined or below threshold. It also rejects early
stop if any task has failed or been blocked. If the predicate is true, the caller marks only
the currently remaining tasks with \texttt{task.optional = True} as skipped. No mandatory task
is skipped by this branch.
\end{proof}

\begin{lemma}[Termination on planner-generated plans]
For every planner-generated plan, \texttt{Scheduler.run} terminates and returns a finite
\texttt{ScheduleResult}.
\end{lemma}

\begin{proof}
By the acyclicity lemma, every nonempty remaining set contains at least one task whose
dependencies are satisfied or one task made blocked by failed predecessors. In each outer-loop
iteration, at least one task is removed from the finite set \texttt{remaining}: either by
blocking, by execution completion/failure, or by early-stop skipping. Since the initial number
of tasks is finite, the loop terminates after finitely many removals.
\end{proof}

\section{Remediation Safety}
\begin{proposition}[Only well-formed proposals reach the executor]
The remediation executor is called only for objects that were successfully materialized as
\texttt{ActionProposal} values.
\end{proposition}

\begin{proof}
The remediation agent invokes the executor only inside the loop
\texttt{for proposal in proposals}. That list is produced only after
\texttt{validate\_action\_proposals(raw)} returns successfully. If parsing or validation does not
produce such proposals, the loop body does not run and the executor is not called.
\end{proof}

\begin{proposition}[Exact authorization condition]
For an \texttt{ActionExecutionRequest} \(q\), \texttt{ExecutionPolicy.decide(q)} yields
\texttt{should\_execute = True} if and only if all four conditions hold:
\begin{enumerate}[label=(\roman*)]
\item execution policy is enabled;
\item the action type is allowlisted;
\item the adapter supports the action type;
\item the request has \texttt{allow\_execution = True}.
\end{enumerate}
\end{proposition}

\begin{proof}
Inspect the decision tree in \texttt{ExecutionPolicy.decide}. The function returns
\texttt{should\_execute = False} immediately whenever any one of the four listed conditions is
false, and returns \texttt{should\_execute = True} only in the final branch reached exactly when
all four are true. Therefore the condition is necessary and sufficient.
\end{proof}

\begin{corollary}[No unauthorized remediation execution attempt]
On the implemented remediation path, no adapter execution is attempted unless the policy is enabled,
the action is allowlisted, the adapter supports it, and auto-execution was explicitly enabled.
\end{corollary}

\begin{proof}
The remediation path is
\[
\texttt{RemediationAgent} \rightarrow \texttt{ActionExecutorTool.run} \rightarrow \texttt{adapter.execute}.
\]
The tool calls \texttt{adapter.execute} only in the branch where the policy decision has
\texttt{should\_execute = True}. By the previous proposition, this happens exactly under the four
authorization conditions. Moreover, if no explicit policy object is supplied, the remediation
agent constructs one with \texttt{allowed\_actions = ()}, so execution still fails closed by
default. Hence no unauthorized adapter execution attempt can occur on this path.
\end{proof}

\begin{remark}
This is a code-level authorization theorem. It does not assert that an authorized external API
call must succeed remotely; it asserts that the code does not even attempt adapter execution
before the authorization predicate is satisfied.
\end{remark}

\section{Memory, Correlation, and Replay}
\begin{lemma}[At most one stored memory per run identifier]
The memory store contains at most one \texttt{IncidentMemory} row per \texttt{run\_id}.
\end{lemma}

\begin{proof}
The \texttt{incident\_memory} schema declares \texttt{run\_id} unique. The write path performs an
\texttt{INSERT ... ON CONFLICT(run\_id) DO UPDATE}. Therefore distinct rows with the same
\texttt{run\_id} cannot coexist.
\end{proof}

\begin{lemma}[Successful memory write is post-checked]
If \texttt{MemoryStore.write\_memory} returns normally with value \(m\), then \(m\) is retrievable
from the store by its \texttt{run\_id}.
\end{lemma}

\begin{proof}
After issuing the insert or upsert, \texttt{write\_memory} immediately calls
\texttt{get\_memory\_by\_run\_id(record.run\_id)}. If that query returns \texttt{None}, the method
raises \texttt{RuntimeError}. Therefore any normal return is accompanied by a successful
postcondition check.
\end{proof}

\begin{lemma}[Baseline monotonicity on the application path]
Under Assumption 3, each completed application-path memory write increments the observed
baseline count by exactly one for each unique extracted entity in that memory.
\end{lemma}

\begin{proof}
The application path obtains entities from \texttt{extract\_entities\_from\_graph}, which emits
unique lists. Then \texttt{MemoryStore.\_update\_baselines} iterates over each entity value once
and either inserts a new observed baseline with incident count \(1\) or updates the existing row
by adding \(1\). Because no duplicate entity value appears within a bucket on the application
path, the increment occurs exactly once per unique extracted entity.
\end{proof}

\begin{lemma}[Correlation returns only matching prior context]
Every incident returned by \texttt{CorrelationService.get\_prior\_context} matches at least one
entity extracted from the queried alert, and no run identifier is returned twice.
\end{lemma}

\begin{proof}
For each extracted entity bucket, the correlation service queries
\texttt{list\_memories\_for\_entity(canonical\_type, value)}. That store method includes a memory
only if the queried value appears in one of the stored entity lists under the same canonical
entity type. The correlation service tracks a \texttt{seen\_runs} set and skips any memory whose
\texttt{run\_id} has already been emitted. Therefore every returned incident matches at least one
queried entity and duplicate runs are eliminated.
\end{proof}

\begin{lemma}[Replay preserves stored alert semantics]
If a run identifier exists in memory, \texttt{replay\_investigation(run\_id, ...)} reconstructs
the stored alert from \texttt{alert\_json} and passes that reconstructed alert to
\texttt{run\_investigation}.
\end{lemma}

\begin{proof}
\texttt{replay\_investigation} first retrieves the memory by \texttt{run\_id}. It then calls
\texttt{\_reconstruct\_alert(memory.alert\_json)} and passes the resulting alert object directly to
\texttt{run\_investigation}. If the run identifier is missing, it raises \texttt{ValueError}
instead of inventing a replacement alert. Therefore replay preserves the persisted alert
semantics or fails explicitly.
\end{proof}

\section{Main Theorem}
\begin{theorem}[Implemented control-plane correctness of \texttt{soc-agent}]
Under Assumptions 1--3, for every call to
\[
\texttt{run\_investigation(config, alert, dry\_run, ...)}
\]
that returns normally, the following statements are true:
\begin{enumerate}[label=(\arabic*)]
\item the task DAG actually executed is a deterministic function of the input alert and not of
      the commander's LLM output;
\item every task executed by the scheduler begins only after its dependencies are completed or
      skipped;
\item the scheduler terminates in finitely many steps on the planner-generated plan and assigns
      each scheduled task exactly one final result among completed, failed, blocked, or skipped;
\item adapter execution is attempted on the remediation path if and only if the explicit
      authorization conditions of the execution policy are satisfied;
\item if memory is enabled, every normal return from \texttt{write\_memory} implies that the run's
      incident memory is persisted and retrievable by \texttt{run\_id};
\item if a replay is later requested for that persisted \texttt{run\_id}, the replay path
      reconstructs the stored alert semantics instead of manufacturing a different alert.
\end{enumerate}
\end{theorem}

\begin{proof}
Item (1) follows from the planner determinism lemma and the commander LLM-independence lemma.
Item (2) follows from the scheduler dependency-safety lemma. Item (3) follows from the blocked
propagation, bounded-retry, early-stop, and termination lemmas. Item (4) follows from the exact
authorization proposition and its corollary. Item (5) follows from the memory post-check lemma.
Item (6) follows from the replay lemma. The statements are composable because the application
pipeline wires these subsystems in series without any contradictory branch that bypasses the
proved guards. Therefore the theorem holds.
\end{proof}

\section{Empirical Validation of the Proof Obligations}
The proofs above are deductive. The current workspace also provides empirical validation that the
implemented code conforms to the proved obligations.

On March 27, 2026, the command
\[
\texttt{pytest soc-agent/tests -q}
\]
completed with
\[
\texttt{225 passed,\ 8 skipped}.
\]

The passing tests include, among others:
\begin{enumerate}[label=(\arabic*)]
\item planner determinism and expected DAG structure;
\item scheduler concurrency, retry bounds, timeout handling, and optional-only early stop;
\item execution-policy gating and adapter routing;
\item memory-store persistence, baseline updates, and correlation behavior;
\item entity extraction and replay reconstruction;
\item approval identity, approval queue, worker queue, and application-level investigation flow.
\end{enumerate}

\begin{remark}
The tests do not replace the proof. They provide implementation evidence that the source under
inspection matches the logical properties proved above.
\end{remark}

\section{Conclusion}
The strongest defensible statement is the following:

\begin{quote}
\texttt{soc-agent} is mathematically and logically correct with respect to its implemented
control-plane semantics: deterministic planning, dependency-safe scheduling, bounded failure
behavior, fail-closed remediation authorization, persistent incident memory, correlation, and replay.
What is \emph{not} proved is universal empirical threat-detection accuracy over arbitrary
real-world incidents.
\end{quote}

That statement is both rigorous and true for the code inspected here.

\end{document}
